\section{Introduction}\label{sec:intro}

A transformer's internal components need not retain the same causal role after the circuit around them is perturbed. Ablating one component can change the contribution of others, producing self-repair, backup behavior, or more general conditional importance \citep{mcgrath2023hydra,mcdougall2023copysuppression,rushing2024selfrepair,gong2026coax}. Detecting such a change is useful, but it leaves a mechanistic question unresolved: \emph{why did the selected component become more or less useful?}

For an attention head, several explanations are possible. Its task-relevant attention score may increase; competing scores may decrease; value content may change; or essentially the same head output may become more or less useful because the surrounding network has changed. These possibilities are not interchangeable. A visible increase in attention can occur without positive functional recruitment, and two heads that satisfy the same conditional-recruitment criterion can in principle do so through different internal changes.

We introduce \textbf{FALI (Factorized Attention-Level Intervention)}, a hypothesis-driven framework for resolving that ambiguity. FALI starts after a candidate attention head and a task-relevant target edge have been identified. It compares an intact state with a perturbed state and decomposes the resulting change in the candidate's \emph{output-grounded causal usefulness}. The candidate's attention computation is factorized into target scores $T$, rival scores $R$, other valid scores $O$, value content $V$, and surrounding-network context $C$. Evaluating the resulting $2^5$ counterfactual states defines a causal game whose interactions are allocated with the standard Owen value. The fine score contributions reconcile exactly with a coarser SCORE/VALUE/CONTEXT decomposition.

The key distinction is between \emph{localizing conditional importance} and \emph{explaining it}. Path and activation patching identify behaviorally relevant components or paths \citep{goldowskydill2023pathpatching,zhang2024activationpatching,heimersheim2024activationpatching}; conditional co-ablation identifies components whose importance changes after a primary intervention \citep{gong2026coax}; interaction-aware mediation analyses show that a component's measured effect can depend strongly on the rest of the model \citep{vaidyanathan2026curse}. FALI is downstream of these approaches: given a selected conditional change, it asks which architecture-aligned changes account for it under an explicit intervention model.

We first validate the intervention semantics in controlled systems. Deterministic synthetic attention models isolate each of $T,R,O,V,C$ and a known cross-factor interaction, while a small randomly initialized transformer checks score-to-pattern reconstruction, value-to-output reconstruction, endpoint recovery, intervention isolation, batch-specific score partitions, and consistency between the fine Owen game and its coarse quotient game. These tests establish that the implementation recovers the factors it manipulates; they do not claim that the factorization is unique or that hybrid states are naturally occurring model states.

We then apply FALI to GPT-2 Small on indirect object identification (IOI) \citep{wang2023ioi}. Corrected functional analysis separates attention adaptation from recruitment: after the canonical primary Name Movers are removed, L10H2 becomes more causally useful, whereas L10H10 increases target attention even more strongly while its causal usefulness decreases. FALI assigns both heads a target-positive/rival-negative score pattern, showing that even a seemingly favorable score reorganization does not itself define recruitment. For positively recruited L10H2, the main fine-score contributions are $\phi_T>0$ and $\phi_R<0$ across the canonical confirmation cohort.

The decisive test is cross-setting transfer. We prospectively test the GPT-2 target-positive/rival-negative mechanism in Pythia-160M induction. Positive conditional recruitment appears again, but the fine score signs reverse: $\phi_T<0$ and $\phi_R>0$. Thus the preregistered mechanism-transfer hypothesis fails. The shared functional label---conditional recruitment---does not imply a shared fine-grained mechanism.

Our contributions are threefold. First, we define an output-grounded counterfactual estimand for the cross-state change in a selected attention head's causal usefulness. Second, we factor that change into score, value, and surrounding-context contributions, with target/rival/other resolution inside SCORE and coarse--fine consistency via standard Owen attribution. Third, we show empirically that two settings with positive conditional recruitment can exhibit opposite target/rival score mechanisms. To our knowledge, this combination provides a distinct mechanism-resolution layer for conditionally important attention heads, rather than a new theory of self-repair, a new circuit-discovery method, or new Shapley/Owen mathematics.
